\begin{minipage}[t]{0.65\textwidth}
Albert a acquis un étang d'une surface de 2\,000 m$^2$.

Le jour de son anniversaire, un dimanche, il installe des nénuphars sur une surface de 200 m$^2$.

\begin{enumerate}
\item Le dimanche d'après, la surface des nénuphars a augmenté de 40 m$^2$.
\begin{enumerate}
\item Quel pourcentage d'augmentation cela représente-t-il ?
\item Quelle est à présent la surface occupée par les nénuphars ?
\end{enumerate}
\end{enumerate}
\end{minipage}
\hfill
\begin{minipage}[t]{0.3\textwidth}
\vspace{0pt}
\begin{center}
\begin{tikzpicture}[scale=0.75]
\draw[fill=cyan!60, draw=black, line width=1pt] (0,0) ellipse (2.5cm and 1.5cm);
\draw[fill=green!70!black, draw=black, line width=1pt] (1.5,0) ellipse (1cm and 0.5cm);
\draw[-{Triangle[length=3mm,width=3mm]}, line width=1.25mm, black] (-1.5,1.6) -- (0,0.75);
\draw[-{Triangle[length=3mm,width=3mm]}, line width=1.25mm, black] (2.2,1.6) -- (1.5,0);
\node[above] at (-1.5,1.6) {Etang};
\node[above] at (2.2,1.6) {Nénuphars};
\end{tikzpicture}
\end{center}
\end{minipage}

\begin{enumerate}[resume]
\item Dans cette question, on suppose que la surface occupée par les nénuphars augmente de 40 m$^2$ chaque semaine, depuis la date de l'anniversaire, tant que cela est possible.
\begin{enumerate}
\item Quelle sera la surface occupée par les nénuphars 10 semaines après l'anniversaire ?
\item Est-il possible qu'un dimanche, la surface occupée par les nénuphars soit égale à 580 m$^2$ ? Justifier.
\item Au bout de combien de semaines, l'étang sera-t-il entièrement recouvert de nénuphars ?
\end{enumerate}

\item Dans cette question, on suppose que la surface occupée par les nénuphars augmente de 20\% chaque semaine, depuis la date de l'anniversaire, tant que cela est possible.
\begin{enumerate}
\item Quelle sera la surface occupée par les nénuphars 2 semaines après l'anniversaire ?
\item On considère un entier naturel $n$. Déterminer, en fonction de $n$, la surface occupée par les nénuphars $n$ semaines après l'anniversaire ?
\item Au bout de combien de semaines, l'étang sera-t-il entièrement recouvert par les nénuphars ? On pourra s'aider du tableau ci-dessous.

\begin{center}
\begin{tabular}{|c|c|c|c|c|c|c|c|c|c|}
\hline
$n = $ & 0 & 1 & 2 & 5 & 10 & 12 & 13 & 14 & 15 \\
\hline
$1,2^n \approx $ & 1 & 1,2 & 1,44 & 2,49 & 6,19 & 8,92 & 10,70 & 12,84 & 15,40 \\
\hline
\end{tabular}
\end{center}
\end{enumerate}

\item Réaliser sur votre copie un \textit{schéma} sur lequel apparaissent l'allure des nuages de points traduisant la progression de la surface occupée par les nénuphars, aussi bien dans le cas de la question 2 que dans le cas de la question 3, et faire figurer le moment où, dans chacun des cas, l'étang est recouvert par les nénuphars.
\end{enumerate}
